% Use
%  pdflatex slides.tex
% to compile this model of slides.
% Use the following first line for animation..
%\documentclass[9pt]{beamer}
% .. or this one for printing
\documentclass[9pt,handout]{beamer}
\usepackage[french]{babel}
\usepackage[utf8]{inputenc}
\usepackage{hyperref}
\usepackage{verbatim}
\usepackage{wasysym}
\usepackage{multicol}
\usepackage{xspace}
\usepackage{listings}
\usepackage{color}
\usepackage{xcolor}
\usepackage{array}
\usepackage{pgf,tikz}
\usepackage{hhline}
\usepackage{colortbl}
\usepackage{rotating}
\usepackage{ragged2e}
\usepackage[font=rmfamily,leftmargin=10pt,rightmargin=10pt]{quoting}

\definecolor{violet}{rgb}{0.5,0,0.5}
\definecolor{lemon}{rgb}{0.9,0.9,0.4}
\definecolor{mygreen}{RGB}{27,156,16}
\definecolor{myred}{RGB}{221,0,17}
\definecolor{fushia}{rgb}{255,0,255}
\definecolor{orangeb}{RGB}{204,85,0}

\DeclareGraphicsExtensions{pdf,eps}

% définition du style Why3
% For listings:
\definecolor{redcomm}{rgb}{0.8,0,0}
\definecolor{byzantium}{rgb}{0.44,0.16,0.39}
\lstdefinelanguage{Why3}
{
 literate=
  {forall}{{$\forall\ $}}1
  {exists}{{$\exists\ $}}1
  {=>}{{$\Rightarrow$}}1
  {==>}{{$\Longrightarrow$}}1
  {>->}{{$\rightarrowtail$}}2
  {<-}{{$\leftarrow\ $}}1
  {<->}{{$\leftrightarrow\ $}}1
  {->}{{$\to\ $}}1
  {<=}{{$\leq\ $}}1
  {>=}{{$\geq\ $}}1
  {<>}{{$\neq\ $}}1
  {\/\\}{{$\wedge\ $}}1
  {|-}{{$\vdash$}}1
  {\\\/}{{$\vee\ $}}1
  {~}{{$\sim$}}1,
 basicstyle=\scriptsize\ttfamily\bfseries,
 sensitive=true,
 morecomment=[l][\color{redcomm}]{//},
 morecomment=[s][\color{redcomm}]{(*}{*)},
 morestring=[b],
 stringstyle=\color{black},
 showstringspaces=false,
 firstnumber=\thelstnumber,
 numberstyle=\small,
 numberblanklines=true,
 showspaces=false,
 showtabs=false,
 xleftmargin=0pt,
 xrightmargin=0pt,
 emph=
 {[1]
  Record,Definition,Inductive,Parameter,Parameters,Axiom,Axioms,Fixpoint,Conjecture,Variable,
  Variables,Hypothesis,Hypotheses,CoInductive,Theorem,Lemma,CoFixpoint,Let,Notation,Eval,Goal,
  Qed,Defined
  },
 emphstyle={[1]\color{red}},
 emph=
 {[2]
  bijective,relations.TotalStrictOrder,numof,module,while,do,begin,done,mutable,ghost,
  set,type,prop,fix,with,cofix,forall,fun,let,if,then,else,match,end,return,as,for,to,in,
  nat,list,int,array,compute,unfold,intro,intros,assert,omega,decompose,try,simpl,or,and,subst,true,false,
  firstorder,nth,val,lemma,use,import,writes,function,axiom
  },
 emphstyle={[2]\color{blue}},
 emph=
 {[3]
  predicate,invariant,variant,requires,ensures,writes
  },
 emphstyle={[3]\color{blue}}
}


\lstnewenvironment{codewhy3}{\lstset{language=Why3,numbers=none}}{\smallskip}

\lstdefinestyle{why3}{language={Why3},%
  numbers=none,
  captionpos=b,
  columns=flexible
  basicstyle=\normalsize\ttfamily\bfseries
}

\lstdefinestyle{why3footnotesize}{language={Why3},%
  numbers=none,
  captionpos=b,
  columns=flexible,
  basicstyle=\footnotesize\ttfamily\bfseries
}

\lstdefinestyle{why3small}{language={Why3},%
  numbers=none,
  captionpos=b,
  columns=flexible,
   basicstyle=\small\ttfamily\bfseries
}

\lstdefinestyle{why3scriptsize}{language={Why3},%
  numbers=none,
  captionpos=b,
  columns=flexible,
   basicstyle=\scriptsize\ttfamily\bfseries
}

\lstdefinestyle{why3tiny}{language={Why3},%
  numbers=none,
  captionpos=b,
  columns=flexible,
   basicstyle=\tiny\ttfamily\bfseries
}


\newcommand{\why}[1]{\lstinline[style=why3small]{#1}}
\newcommand{\whyf}[1]{\lstinline[style=why3footnotesize]{#1}}

\bibliographystyle{alpha}
\setbeamertemplate{bibliography item}[text]
\renewcommand\newblock{\hskip .11em\@plus.33em\@minus.07em}

\title[\textit{Changer le titre court ici}]{\textit{Changer le titre long ici}}

\author[P. Nom]{\textit{Prénom1 Nom1 et Prénom2 Nom2 etc}\\
\textit{Groupe TP}\\
\vspace{4mm}
UMLP, Besan\c{c}on,
France}

\date{\textit{\today}
\begin{center}
\textit{(Toutes les parties en italique, dont celle-ci, sont à remplacer par vos propres textes)}
\end{center}
}

\begin{document}

\maketitle

% Plan de la présentation

\setcounter{tocdepth}{1}     % Profondeur de la table des matieres

\begin{frame}
\frametitle{Plan}
 \tableofcontents
\end{frame}


\section[Contexte]{Contexte (\textit{de 1 à 4 slides})}

\begin{frame}
\frametitle{\secname}

\textit{Décrire ici le contexte scientifique du sujet choisi, avec l'apport des références principales}

\bigskip


\vspace{40pt}

\begin{flushright}
\textit{Prénom Nom de l'auteur-présentateur du slide, à mettre sur chaque slide}
\end{flushright}
\end{frame}


\section{Existant/Définition/Vocabulaire (\textit{de 0 à 4 slides})}

% Exemple de slide qui cite une référence
\begin{frame}[fragile]
\frametitle{\secname}
\framesubtitle{\textit{exemple de sous titre, à remplacer} : Une permutation est une bijection}

\textit{Selon le sujet, on peut avoir besoin d'un ou plusieurs slides comme celui-ci, introduisant des notions, définitions, notations, références, etc. Remplacer ce modèle selon les besoins de votre sujet, avec au moins une référence.}

\justify \begin{quoting}[leftmargin=6pt,rightmargin=6pt] ``the signature used in
\textsf{TOOB} (\textsf{T}heory \textsf{O}f \textsf{O}ne \textsf{B}ijection)
consists of a single [\ldots] relation between two elements [that] indicates
that the first one is sent to the second one by the permutation viewed as a
bijection.''~\cite[partie 2.1]{ABF20}
\end{quoting}

\pause

\begin{center}
\begin{tabular}{|c|c|}
\hline
Théorie d'une relation & Modèle $R^{f}$ associé à une bijection $f$ \\
binaire bijective $R$  & $x~R^{f} y \Leftrightarrow f(x) = y$

\\
\hline
\begin{lstlisting}[style=why3footnotesize]
module TOOB
  type t
  predicate rel t t
  clone export Bijectivity with
    type t = t, type u = t,
    predicate rel = rel, axiom .
end
\end{lstlisting}
&
\begin{tabular}{l}
\noindent
\begin{lstlisting}[style=why3footnotesize]
type t
val constant f : map t t
axiom f_bij : bijective f
predicate rel_f (x y: t) = (f[x] = y)
clone TOOB with type t = t, 
  predicate rel = rel_f
\end{lstlisting}
\end{tabular}
\\
\hline
\end{tabular}
\end{center}
\pause


\noindent\rule{\textwidth}{0.4pt}

% Copier ci-dessous les blocs \bibitem générés dans le fichier .bbl par la
% compilation de article.tex avec pdflatex
\begin{scriptsize}
\begin{thebibliography}{ABF20}

\bibitem[ABF20]{ABF20}
Michael Albert, Mathilde Bouvel, and Valentin Féray.
%\newblock 
Two first-order logics of permutations.
%\newblock 
{\em Journal of Combinatorial Theory, Series A}, 171:105158, April
  2020.
\newblock \url{https://arxiv.org/abs/1808.05459v2}.
\end{thebibliography}
\end{scriptsize}
\end{frame}

\section{Problématique (\textit{1 seul slide})}

\begin{frame}
\frametitle{\secname}

\textit{Présenter ici vos questions de recherche}

\end{frame}


\section{\textit{Traitement du problème et programme de recherche} (\textit{de 1 à 4 slides})}

\begin{frame}
\frametitle{\secname}

\textit{Présenter ici votre compréhension du problème, vos contributions, vos codes produits, vos formalisations éventuelles en Why3 et/ou Prolog, puis proposer un programme de recherche qui permettrait de prolonger cette étude lors des prochains semestres du master. Vous pouvez aussi suggérer des besoins de formation complémentaire, de lectures complémentaires à faire, de notions à acquérir, d'expérimentations à réaliser, de discussions à mener avec des chercheurs, etc}

\end{frame}

\section{Synthèse/Conclusion (\textit{1 seul slide})}

\begin{frame}
\frametitle{\secname}

\textit{Comparer ici le travail accompli avec les questions de recherche initiales, afin de mesurer l'avancement du travail et d'évaluer la faisabilité/difficulté du reste.}

\bigskip

Conclusions personnelles

\textit{Ensuite, par un item par membre de l'équipe, résumer le rôle, les apports et les difficultés de chacun}
\begin{itemize}
\item \textit{Prénom1 Nom1 : ..}
\item \textit{Prénom2 Nom2 : ..}
\item \textit{etc}
\end{itemize}
\end{frame}

\end{document}
